% =============================================================
\section{Background and related work}\label{sec:background}
% =============================================================

\subsection{Extracellular potential from point-source electrodes}

At neural stimulation frequencies, displacement currents in tissue are
negligible compared with conduction currents, and Maxwell's equations
reduce to a quasi-static formulation in which the extracellular potential
satisfies Laplace's equation away from sources. Modelling each electrode
contact as a point source in an infinite, homogeneous, isotropic
conductor of conductivity $\sigma$ gives the closed-form monopole
$\phi(\mathbf{x}) = I / (4\pi\sigma\,d)$, with $d$ the source-to-field
distance~\cite{10.1093/acprof:oso/9780195058239.001.0001,10.1093/acprof:oso/9780195050387.001.0001}.
Multiple electrodes follow by superposition. We use this model throughout; the homogeneity and isotropy
assumptions are limitations we return to in \cref{sec:limitations}.

\subsection{The activating function}

For a uniform straight cable in an extracellular field $\phi(x,t)$,
Rattay~\cite{rattay1986analysis} showed that the driving term for
transmembrane voltage is proportional to
$\partial^2 \phi / \partial x^2$ along the cable; McNeal's earlier
compartmental treatment of myelinated axons under monopolar
stimulation~\cite{Mcneal1976AnalysisOA} contains the same idea in
discrete form. On an arbitrary compartmental tree the discrete
analogue is $[G\boldsymbol{\phi}]_j$, where $G$ is the axial diffusion
operator that the cable solver already maintains; the continuous second
derivative is the special case of a uniform unbranched cable. We exploit
this in \cref{sec:method}: rather than recomputing
$\partial^2 \phi / \partial x^2$ along each branch, we apply Jaxley's
$G$ directly to the sampled extracellular potential.

\subsection{Differentiable neuron simulators}

Jaxley~\cite{deistler2025jaxley} is the only published JAX-native
compartmental simulator we are aware of that supports HH-style channel
kinetics, branched morphologies, and reverse-mode autodiff through the
time integrator. Diffrax~\cite{DBLP:journals/corr/abs-2202-02435} provides
differentiable ODE solvers in JAX, but it does not model compartmental
neurons. Brian2~\cite{Stimberg2019} is
Python-native and considerably faster than NEURON, but it is not built
for automatic differentiation. We reuse Jaxley's solver through its public
stimulus API rather than forking the simulator: the extracellular extension
is a JAX module that composes with the rest of the framework without
modifying its internals.

